\section{Background and Related Works}

\subsection{Missing Data in Medical and Longitudinal Imaging}

Missing data is a common and significant challenge in clinical imaging studies, particularly in longitudinal settings where subjects undergo repeated scans over time. Data may be missing at one or more timepoints due to various factors such as motion artifacts, missed visits, or participant dropout. This problem is especially pronounced in aging or disease-prone populations, where frequent imaging is required to monitor disease progression, thereby increasing the likelihood of incomplete scan sequences. 

Traditional approaches to handling missing data, such as list wise deletion where subjects with incomplete imaging records are excluded, can substantially reduce the effective sample size and introduce selection bias, ultimately compromising the validity and generalizability of the study findings~\cite{}. Alternative imputation techniques, including last observation carried forward and linear interpolation, have been applied in certain contexts. However, these methods are primarily designed for structured, low-dimensional data and are not well suited to the spatial and contextual complexity inherent in high-dimensional modalities such as three-dimensional volumetric imaging.

In medical imaging, missing data can arise at multiple levels, including within individual scans, for example when slices are absent or corrupted, and across temporal sequences, such as when entire follow-up volumes are unavailable. While deep learning models such as conditional generative adversarial networks and convolutional neural networks have demonstrated success in intra-scan slice imputation, these approaches do not readily extend to the more complex task of reconstructing full three-dimensional scans at missing timepoints in a longitudinal imaging sequence.

% Recent models have attempted to address this challenge of full-scan imputation in longitudinal imaging. ReMiND~\cite{} uses a diffusion-based approach to generate intermediate MRI volumes conditioned on neighboring timepoints. However, it operates in intensity space and does not explicitly enforce anatomical or temporal consistency. NODER~\cite{} takes a different approach by modeling trajectories in a latent space per subject and regressing the missing scan directly, again focusing on image-level synthesis without modeling anatomical transformations. In contrast, our proposed method, SynBADD, addresses scan imputation in deformation space by predicting stationary velocity fields (SVFs) that warp a baseline scan to the target timepoint. This allows anatomically consistent interpolation across discrete timepoints using both past and future context, offering a structure-aware alternative to intensity-based synthesis.



\subsection{Predicting Future Scans in Longitudinal Medical Imaging}

% Predictive modeling of future scans enables proactive assessment of disease progression. Early approaches employed parametric methods such as linear mixed-effects models~\cite{} and geodesic regression~\cite{} to model trajectories of derived features like brain volume or cortical thickness. While interpretable, these methods often assume smooth, low-dimensional trends that may not capture the heterogeneous, nonlinear nature of neurodegenerative processes. More recently, deep learning models have been introduced to directly predict future scans from imaging data. Architectures based on recurrent networks and convolutional encoder-decoders have shown success in extrapolating anatomical progression from earlier timepoints~\cite{}. However, many of these models are trained in a unidirectional manner, relying solely on past scans for forward prediction, which limits their applicability when intermediate timepoints are missing. 

Predictive modeling of future scans enables proactive assessment of anatomical and pathological progression. Early approaches in longitudinal imaging relied on parametric models such as linear mixed-effects models~\cite{} and geodesic regression~\cite{} to capture temporal trends in derived clinical or morphometric features, such as brain volume or cortical thickness. While interpretable, these methods often assume smooth, low-dimensional trajectories that may not reflect the heterogeneous and nonlinear progression patterns typical of neurodegenerative diseases. More recent deep learning approaches directly predict future scans from past imaging data, using architectures such as recurrent neural networks and convolutional encoder-decoders~\cite{}. However, many of these models are trained in a unidirectional manner, relying solely on preceding scans for forward prediction, which limits their applicability when intermediate timepoints are missing.

Recent models such as ReMiND~\cite{} and NoDeR~\cite{} address this limitation by incorporating temporal context for scan imputation. ReMiND employs a diffusion model conditioned only on the immediate past and future scans, synthesizing missing images in intensity space without enforcing anatomical correspondence. NODER, in contrast, learns a subject-specific trajectory in the latent space of stationary velocity fields (SVFs), enabling intermediate scan prediction through curve fitting. However, its coarse supervision may lead to oversmoothing or inaccurate estimation of anatomical changes. SynBADD adopts a distinct approach by learning deformation fields—parameterized as stationary velocity fields (SVFs)—from the baseline to the missing timepoint, while leveraging the full longitudinal context. By operating in deformation space rather than directly synthesizing intensities, SynBADD preserves subject-specific anatomy, maintains diffeomorphic consistency, and supports temporally coherent interpolation at discrete timepoints.



% Recent models such as ReMiND~\cite{} and NoDeR~\cite{} attempt to address this by incorporating temporal context. ReMiND employs a diffusion model to synthesize missing scans based on a pair of neighboring scans, one before and one after the missing timepoint. However, it operates in intensity space and lacks explicit constraints to preserve anatomical correspondence across time. NoDeR, on the other hand, learns a subject-specific latent trajectory by fitting a continuous curve through stationary velocity fields (SVFs). While this provides a mechanism for interpolation, it may over-smooth anatomical changes or overestimate the trajectory due to limited supervision. SynBADD differs in both formulation and design. It models missing anatomy by learning deformation fields—parameterized as SVFs—that transform the baseline scan to the target timepoint. Unlike ReMiND, which conditions on only two neighboring scans, SynBADD leverages the full longitudinal context and supports interpolation at arbitrary discrete timepoints. By operating in deformation space rather than image space, SynBADD preserves subject-specific anatomy, enforces diffeomorphic structure, and yields temporally consistent interpolations aligned with realistic patterns of anatomical change.





% Recent methods such as ReMiND~\cite{} and NoDeR~\cite{} have addressed intermediate scan imputation by incorporating temporal context. ReMiND uses a diffusion model to generate a missing scan conditioned on neighboring scans. However, as it synthesizes images directly in intensity space from scratch, it does not explicitly enforce anatomical consistency. NoDeR, in contrast, performs subject-specific regression in a latent space of stationary velocity fields (SVFs), allowing intermediate scans to be predicted along a learned latent trajectory. 

% SynBADD adopts a different formulation and objective by modeling missing anatomy through anatomically plausible deformation fields, parameterized as SVFs, between observed scans. Unlike ReMiND, which conditions only on immediate temporal neighbors, SynBADD leverages the full longitudinal context to enable interpolation at multiple discrete timepoints. By operating in deformation space, it preserves subject-specific anatomical consistency, maintains diffeomorphic structure, and supports temporally coherent synthesis across the entire timeline.




\subsection{Generative Models for Medical Image Synthesis}

Generative models are increasingly being integrated into medical image analysis workflows, supporting tasks such as modality translation, super-resolution, anomaly detection, and missing data imputation. Common generative frameworks include Variational Autoencoders (VAEs), Generative Adversarial Networks (GANs), and diffusion models. VAEs~\cite{} provide latent representations that support interpolation and uncertainty modeling, but their outputs are often blurry due to the regularization imposed by their probabilistic formulation. Despite this, VAEs have been used effectively in tasks involving anatomical variation and longitudinal modeling~\cite{}. GANs~\cite{}, in contrast, generate sharper images through adversarial training and have been applied to tasks such as cross-modality synthesis~\cite{}, slice inpainting~\cite{}, and super-resolution~\cite{}. However, they are prone to challenges such as mode collapse and training instability, which can limit their reliability and generalizability in clinical applications. Diffusion models~\cite{} have recently gained traction for their strong performance in generating high-resolution images. These models synthesize data by iteratively denoising random noise and are known for their stable training dynamics and flexible conditioning mechanisms. However, their applicability can be limited by slow generation times, primarily due to the computational cost of iterative reverse sampling, which becomes even more pronounced when generating large three-dimensional volumetric scans. ReMiND~\cite{} applies diffusion modeling to the task of longitudinal imputation by synthesizing intermediate MRI slices in two dimensions, which are stacked along the depth axis in groups of 12 slices to approximate a three-dimensional volume. The generation is conditioned on both the preceding and succeeding timepoint scans. However, it operates in image space and is designed to interpolate between a fixed pair of scans, rather than leveraging the full longitudinal sequence for context-aware imputation. Moreover, it lacks explicit constraints to ensure anatomical coherence across the full 3D volume or over time. In longitudinal imaging, preserving subject identity and modeling smooth anatomical progression is essential—especially in neurodegenerative studies where changes are subtle and individualized. Direct voxel-level synthesis may overlook underlying structural transformations. SynBADD addresses this by shifting the generative focus from intensities to deformations. It learns to produce stationary velocity fields (SVFs) that deform a baseline scan to approximate a missing timepoint, ensuring diffeomorphic consistency and temporally structured interpolation across multiple, discretized timepoints.



\subsection{Deformation-Based Approaches in Image Modeling}

Deformation-based models describe anatomical variation by estimating spatial transformations between images. Classical frameworks such as Large Deformation Diffeomorphic Metric Mapping (LDDMM)~\cite{} model smooth, invertible deformations that preserve topological structure, making them well-suited for tasks like registration and morphometric analysis. More recently, neural network-based approaches like VoxelMorph~\cite{} have enabled fast, learning-based registration by training convolutional networks to predict deformation fields directly from image pairs. These methods significantly reduce inference time and support population-level learning, but are typically designed for pairwise alignment and do not naturally generalize to synthesis or imputation tasks. In this work, we extend deformation modeling from alignment to generation. Rather than predicting a deformation between two observed scans, we train a model to synthesize the deformation that maps a baseline image to a missing timepoint. This requires deformations that are both anatomically meaningful and temporally consistent. To this end, we represent transformations using stationary velocity fields (SVFs), which define time-invariant velocity fields whose exponential map yields smooth, diffeomorphic transformations. Compared to displacement fields, SVFs offer a compact, invertible, and topologically stable representation of anatomical change. Crucially, our model predicts deformations from the baseline to each target timepoint, rather than estimating them piecewise between adjacent scans. This baseline-referenced formulation maintains a consistent anatomical coordinate frame across all predictions, improving interpretability and avoiding the accumulation of errors that can arise from chaining local transformations. Unlike sequential methods, which may over-smooth progression or drift in the presence of irregular sampling, our approach enforces a single, well-defined transformation per timepoint. This results in anatomically coherent and temporally aligned interpolation, making it well-suited for modeling longitudinal trajectories in neuroimaging.




% Deformation-based models describe anatomical variation by estimating spatial transformations between images. Classical frameworks such as Large Deformation Diffeomorphic Metric Mapping (LDDMM)~\cite{} model smooth, invertible deformations that preserve topological structure, making them well-suited for tasks like registration and morphometric analysis. More recently, neural network-based approaches like VoxelMorph~\cite{} have enabled fast, learning-based registration by training convolutional networks to predict deformation fields directly from image pairs. These methods significantly reduce inference time and support population-level learning, but are typically designed for pairwise alignment and do not generalize naturally to synthesis or imputation tasks. In our work, we extend deformation modeling from an alignment problem to a generative one. Rather than predicting a deformation between two given scans, we train a model to synthesize the deformation that would map a baseline image to a missing timepoint. This requires the deformation to be both anatomically meaningful and temporally consistent. To this end, we represent deformations using stationary velocity fields (SVFs). 
% SVFs define a stationary (time-invariant) velocity field whose exponential map yields a smooth, diffeomorphic transformation. Compared to displacement fields, SVFs provide a compact and topologically stable representation of anatomical change, ensuring smoothness, invertibility, and structural consistency. In our formulation, deformations are predicted from the baseline scan to the target timepoint, rather than estimated piecewise between adjacent scans. This baseline-referenced approach maintains a consistent anatomical coordinate frame across all predictions, improving interpretability and avoiding error accumulation from chaining multiple local transformations. Unlike methods that rely on sequential fitting of intermediate deformations, which may over-smooth progression or suffer from drift in irregularly sampled data, our strategy enforces a single, well-defined transformation per timepoint. This results in anatomically coherent and temporally aligned interpolation suited for longitudinal neuroimaging applications.



% SVFs define a time-invariant velocity field whose exponential map yields a diffeomorphic transformation. Compared to displacement fields, SVFs guarantee smoothness, invertibility, and better preservation of anatomical topology. They also offer a compact and physically grounded parameterization of deformation, which is advantageous for modeling longitudinal progression in brain anatomy. By predicting SVFs instead of voxel intensities, our method, SynBADD, ensures that interpolated scans remain anatomically coherent and structurally faithful to the subject-specific evolution. This deformation-centric formulation provides a principled way to model missing data in longitudinal neuroimaging.



\subsection{Alzheimer's Disease and Longitudinal Imaging}

Alzheimer’s disease (AD) is a progressive neurodegenerative condition characterized by gradual, region-specific brain atrophy—most notably affecting the hippocampus, entorhinal cortex, along with certain cortical structures~\cite{}. These anatomical changes unfold over years and vary considerably across individuals, making longitudinal imaging a critical tool for studying disease progression, stratifying risk, and informing clinical decision-making. The Alzheimer's Disease Neuroimaging Initiative (ADNI) has played a central role in enabling such research by providing a large-scale, multi-site repository of structural MRI scans acquired at multiple timepoints along the disease spectrum. However, as with many longitudinal studies, ADNI is affected by missing data due to participant dropout, missed visits, and irregular scan intervals. The presence of missing data complicates temporal modeling and can introduce bias when models require complete sequences have to rely on naive imputation techniques. In this context, models capable of anatomically consistent imputation are particularly valuable. SynBADD has been designed with this setting in mind, reconstructing plausible intermediate scans by predicting subject-specific deformation fields based on baseline anatomy and available surrounding timepoints. Rather than generating missing scans directly in intensity space, SynBADD models anatomical progression through predicted deformation fields. By learning stationary velocity fields (SVFs) that map the baseline anatomy to a missing timepoint, it captures subject-specific patterns of neurodegeneration in a temporally consistent and anatomically meaningful way. This approach allows SynBADD to impute intermediate scans even when only partial longitudinal data are available, reducing reliance on fully sampled trajectories. The deformation-based formulation preserves structural continuity and supports robust modeling of disease progression, enabling a more reliable use of incomplete longitudinal imaging data sets in downstream clinical analysis.


% Rather than synthesizing images directly, it models how anatomy changes over time by capturing disease-relevant deformation patterns even in the presence of incomplete data. This deformation-based formulation enables robust longitudinal modeling and supports downstream clinical tasks without requiring fully sampled scan trajectories.




% \section{Background and Related Works}

% %%%%%%%%%%%%%%%%%%%%%%%%%%%%%%%%%%%%%%%%%%%%%%%%%%%
% % MISSING DATA
% %%%%%%%%%%%%%%%%%%%%%%%%%%%%%%%%%%%%%%%%%%%%%%%%%%%

% \subsection{Missing Data in Medical and Longitudinal Imaging}

% Missing data is a persistent issue in medical imaging. Scans can be lost due to motion artifacts, missed visits, patient dropout, or equipment failures. A common handling strategy is to simply exclude subjects with incomplete imaging, but this reduces sample size, introduces bias, and wastes valuable data. Traditional imputation approaches, such as last observation carried forward, linear interpolation or K-nearest neighbours, have been applied to missing measurements. However, most such efforts focused on imputation of derived features or summary phenotypes rather than reconstructing the full images. In medical imaging, missing data can occur both within a scan (e.g., missing slices or corrupted regions) and across scans in time (e.g., missing entire follow-up visits). Recent work has explored learning-based methods to address intra-scan gaps.  For example, conditional GANs have been used to synthesize missing MRI slices in a 3D volume by conditioning on neighboring slices and slice position. Such slice-level imputation showed that deep learning can plausibly recover intra-scan information. Nevertheless, imputing an entire 3D scan at a missing time point is far more challenging due to the high dimensionality and complex anatomical variations involved. As longitudinal studies increasingly include aging or disease-prone populations, who are more likely to miss visits, there is a growing need for models that can impute missing scans in a realistic and anatomically faithful manner.


% %%%%%%%%%%%%%%%%%%%%%%%%%%%%%%%%%%%%%%%%%%%%%%%%%%%
% % PREDICTING FUTURE SCANS
% %%%%%%%%%%%%%%%%%%%%%%%%%%%%%%%%%%%%%%%%%%%%%%%%%%%

% \subsection{Predicting Future Scans in Longitudinal Medical Imaging}

% Predictive modeling in longitudinal imaging aims to forecast future anatomical states based on historical imaging data. Early approaches used statistical models such as linear mixed-effects models and geodesic regression to capture temporal trends. While these models provided valuable insights, they often assumed linear progression and lacked to flexibility to model complex, nonlinear anatomical changes.

% Dee

% However, many existing deep learning models focus on forward prediction, estimating future scans based solely on past data. This unidirectional approach may not fully leverage the available information, especially when intermediate scans are missing.


% %%%%%%%%%%%%%%%%%%%%%%%%%%%%%%%%%%%%%%%%%%%%%%%%%%%
% % GENERATIVE MODELS
% %%%%%%%%%%%%%%%%%%%%%%%%%%%%%%%%%%%%%%%%%%%%%%%%%%%

% \subsection{Generative Models for Medical Image Synthesis}

% Variational autoencoders (VAEs) learn a compact latent space but often produce blurry images due to their probabilistic nature. GANs improve realism by using adversarial training and have been widely applied for tasks such as modality translation and image inpainting. Conditional GANs have shown success in filling missing slices or generating modality-specific images (e.g., predicting T2 from T1-weighted MRI). However, GANs suffer from mode collapse and unstable training. Recently, diffusion models have shown superior image quality and stability. These models refine random noise into structured outputs through a series of denoising steps. They offer flexible conditioning and have achieved state-of-the-art results in brain MRI synthesis. ReMiND is an example that uses a diffusion model to impute missing MRIs slice-by-slice. It generates high-fidelity images conditioned on available scans but works in intensity works in intensity space, often without enforcing anatomical or temporal coherence. SynBADD builds on this progress by introducing a deformation-based generative approach. Rather than synthesizing intensities, it generates deformation fields constrained by anatomy, allowing the model to produce realistic, temporally consistent outputs while preserving subject identity.


% %%%%%%%%%%%%%%%%%%%%%%%%%%%%%%%%%%%%%%%%%%%%%%%%%%%
% % DEFORMATION APPROACHES
% %%%%%%%%%%%%%%%%%%%%%%%%%%%%%%%%%%%%%%%%%%%%%%%%%%%

% \subsection{Deformation-Based Approaches in Image Modelling}

% Deformation-based models offer a compelling alternative for modeling anatomical changes by focusing on the transformation between images rather than directing generating the intensities. Classical methods like Large Deformation Diffeomorphic Metric Mapping (LDDMM) model anatomical chnages as smooth, invertible transformations, preserving the topology  and structure of anatomical regions. These methods have been effective in tasks like image registration and morphometric analysis. Such approaches have also been integrated into neural network architectures. [...]

% However, many existing deformation-based models are designed for pairwise image registration and may not effectively handle the imputation of missing intermediate scans. 


% %%%%%%%%%%%%%%%%%%%%%%%%%%%%%%%%%%%%%%%%%%%%%%%%%%%
% % ALZEHEIMER'S DISEASE
% %%%%%%%%%%%%%%%%%%%%%%%%%%%%%%%%%%%%%%%%%%%%%%%%%%%

% \subsection{Alzheimer’s Disease and the Role of Longitudinal Imaging}

% Alzheimer's disease is a progressive neurodegenerative disorder characterzed by structural brain changes, such as hippocampal atrophy and cortical thinning. Longitudinal MRI studies are crucial for tracking these changes and understanding disease progression. However, missing scans in longitudinal AD studies are common and can hinder the accurate modeling of disease trajectories.
